% =====================================================================
%  1bat-ud2-15.tex  —  HORA 15: Taller de preparació de la prova
%  (sense preàmbul: s'inclou des de main.tex)
% =====================================================================
\horatitol{Sessió 15 — Taller de preparació de la prova}%
{Conèixer per endavant l'estructura i els criteris de correcció de la prova, fer un simulacre i traçar un pla d'estudi.}

\subsection{Idea clau}

A les sessions 1–14 hem treballat els quatre blocs de la unitat: llenguatge
algebraic, valor numèric, operacions amb polinomis i arrels. Avui, abans de
la prova de la sessió 16, en coneixem per endavant l'estructura exacta i
\textbf{com es corregirà}, fem un simulacre breu per blocs i acabem amb un
pla d'estudi personal.

\begin{idea}
\textbf{Estructura de la prova (sessió 16) i rúbrica simplificada.}
La prova té quatre blocs, un per criteri, amb el mateix pes:
\begin{itemize}[leftmargin=1.4em, itemsep=2pt]
  \item \textbf{C1 — Llenguatge:} traduir un enunciat a una expressió o
        un polinomi.
  \item \textbf{C2 — Valor numèric:} substituir i calcular respectant la
        jerarquia d'operacions.
  \item \textbf{C3 — Operacions:} sumar, restar o multiplicar polinomis i
        simplificar el resultat.
  \item \textbf{C4 — Arrels:} trobar el llindar d'una prestació i
        interpretar-lo en context.
\end{itemize}
En cada bloc es valoren quatre nivells: \textbf{NA} (no assolit),
\textbf{AS} (assoliment satisfactori), \textbf{AN} (assoliment notable) i
\textbf{AE} (assoliment excel·lent). El nivell més alt \textbf{no} depèn
només d'encertar el resultat: cal justificar cada pas, controlar el signe
i el grau, i interpretar la resposta en el context del problema.
\end{idea}

\subsection{Exemples}

\begin{exemple}[una resposta de nivell AE — completa i justificada]
Enunciat: una prestació es modelitza amb $H(x)=-10x+350$, on $x$ són els
ingressos, en desenes d'euros. Troba el llindar d'exclusió i interpreta'l.

\medskip
\noindent\textbf{Resposta de nivell AE:}
\[
-10x+350=0 \;\Longrightarrow\; -10x=-350 \;\Longrightarrow\; x=35.
\]
El llindar d'exclusió és $x=35$ (desenes d'euros, és a dir, $350\eur$
d'ingressos). Per sota d'aquest valor encara hi ha dret a la prestació;
a partir d'aquest valor, ja no.

Aquesta resposta mostra el \textbf{procediment} (igualar a zero i
aïllar), el \textbf{resultat} ($x=35$) i la \textbf{interpretació} en
context: és exactament el que demana un nivell AE.
\end{exemple}

\begin{exemple}[la mateixa resposta a nivell AS — correcta però incompleta]
Per al mateix enunciat, un alumne escriu únicament:
\[
x=35.
\]
El resultat numèric és \textbf{correcte}, però no s'hi veu el
procediment (com s'ha arribat a $35$?) ni la interpretació (què vol dir
$35$ en el context de la prestació?). Una resposta així només arriba a
nivell \textbf{AS}: el valor final és necessari, però no suficient.
\end{exemple}

\subsection{Exercicis resolts}

\begin{exresolt}[un ítem mixt: del llenguatge al valor numèric]
Un casal de barri cobra $15\eur$ d'inscripció per infant més $3\eur$ per
cada taller addicional. (a) Tradueix el cost total per a $x$ tallers a
una expressió algebraica. (b) Calcula el cost per a una família que
apunta el seu fill a $4$ tallers.
\solucio
(a) La quota fixa és $15\eur$ i el cost variable és $3\eur$ per taller,
de manera que:
\[
C(x) = 15 + 3x.
\]
(b) Substituïm $x=4$:
\[
C(4) = 15+3\per4 = 15+12 = 27.
\]
\resposta{(a) $C(x)=15+3x$. (b) El cost és $27\eur$.}
\end{exresolt}

\begin{exresolt}[un ítem mixt: operacions i una arrel]
El balanç net de dos serveis és $A(x)=x^2-3x-4$ i $B(x)=-x^2+7x-16$, on
$x$ és el nombre de persones usuàries. (a) Calcula el balanç conjunt
$A(x)+B(x)$. (b) Troba el llindar (arrel) del balanç conjunt i
interpreta el resultat.
\solucio
(a) Sumem terme a terme:
\begin{align*}
A(x)+B(x) &= (x^2-3x-4)+(-x^2+7x-16) \\
          &= (1-1)x^2+(-3+7)x+(-4-16) \\
          &= 4x-20.
\end{align*}
El terme de segon grau es cancel·la: el balanç conjunt és un polinomi de
grau~1.

\noindent(b) Igualem a zero per trobar el llindar:
\[
4x-20=0 \;\Longrightarrow\; x=5.
\]
\resposta{(a) $A(x)+B(x)=4x-20$. (b) Amb $5$ persones usuàries el balanç
és exactament $0$; per sota de $5$ el balanç és negatiu (dèficit) i per
sobre, positiu (superàvit).}
\end{exresolt}

\subsection{Exercicis per fer}

\noindent\textit{Simulacre de prova (35 minuts en total): resol, en
parella, un ítem representatiu de cada bloc amb el temps orientatiu
indicat. Verbalitza l'estratègia abans d'escriure la resposta.}

\begin{exercicis}[Simulacre per blocs]
\begin{exlist}
  \item \textbf{Bloc C1 — Llenguatge algebraic (orientatiu: 7 min).} Un
        servei de menjador comunitari cobra $4\eur$ per àpat més una
        quota fixa mensual de $60\eur$. Tradueix el cost total per a $x$
        àpats a una expressió algebraica.

  \item \textbf{Bloc C2 — Valor numèric (orientatiu: 8 min).} El
        copagament d'un servei es calcula amb
        $C(x)=\num{0,05}x^2-10x+500$ (la mateixa fórmula de les sessions
        4 i 11), on $x$ és la renda mensual en euros. Calcula $C(120)$.

  \item \textbf{Bloc C3 — Operacions amb polinomis (orientatiu: 9
        min).} Dos serveis comparteixen recursos. El cost net d'un és
        $A(x)=3x^2-4x+10$ i el de l'altre $B(x)=x^2+6x-5$, on $x$ és el
        nombre de persones ateses. Calcula i simplifica $A(x)-B(x)$.

  \item \textbf{Bloc C4 — Arrels (orientatiu: 11 min).} El cost net
        d'una activitat es modelitza amb $N(x)=x^2-9x+14$, on $x$ és el
        nombre de participants. Troba les arrels amb la fórmula general
        i interpreta el resultat en context.

\textbf{Atenció a la diversitat:} l'alumnat amb més dificultats pot
disposar, durant el simulacre, de la rúbrica d'avui i d'un
formulari-resum (jerarquia d'operacions, regla del signe en restar
polinomis, fórmula de segon grau), retirant-lo progressivament de cara a
la prova real segons el cas.
\end{exlist}
\end{exercicis}

\subsection*{Pla d'estudi personal}

\noindent A partir d'avui, de les graelles d'autoavaluació de les
sessions 5 i 9 i del simulacre que acabes de fer, marca amb una creu el
teu nivell en cada bloc de la unitat.

\begin{center}
\renewcommand{\arraystretch}{1.4}
\begin{tabular}{|p{7.2cm}|c|c|c|}
\hline
\textbf{Bloc} & \textbf{Domino} & \textbf{Quasi} & \textbf{Em costa} \\
\hline
C1 — Llenguatge algebraic & $\square$ & $\square$ & $\square$ \\
\hline
C2 — Valor numèric & $\square$ & $\square$ & $\square$ \\
\hline
C3 — Operacions amb polinomis & $\square$ & $\square$ & $\square$ \\
\hline
C4 — Arrels & $\square$ & $\square$ & $\square$ \\
\hline
\end{tabular}
\end{center}

\noindent Anota els dos o tres tipus d'exercici que repassaràs
especialment abans de la prova:

\medskip
\noindent 1.\ \dotfill

\noindent 2.\ \dotfill

\noindent 3.\ \dotfill

% ---------------------------------------------------------------------
\fullsolucions{Sessió 15}
\begin{solucions}
\begin{sollist}
  \item $C(x)=4x+60$ (en euros, per a $x$ àpats).
  \item $C(120)=\num{0,05}\per\num{14400}-10\per120+500=720-1200+500=20$:
        el copagament és $20\eur$.
  \item $A(x)-B(x)=(3-1)x^2+(-4-6)x+(10+5)=2x^2-10x+15$.
  \item $\Delta=(-9)^2-4\per1\per14=81-56=25$, $\sqrt{\Delta}=5$;
        $x=\dfrac{9\pm5}{2} \Rightarrow x_1=7$, $x_2=2$. El cost net
        s'anul·la amb $2$ o amb $7$ participants; entremig (per exemple,
        $x=4$) el cost net és negatiu (superàvit), i per sota de $2$ o
        per sobre de $7$, positiu.
\end{sollist}
\end{solucions}
